\section*{How to use this book}

The exam is a conversation of about one hour. First your talk and questions on it, then Machine
Learning and Perception, then AI for Simulation. The examiner picks a topic and follows your answer
with the next question.

\textsf{Every answer has the same five parts.}
\begin{itemize}
\item \textsf{\color{qcol}Core} -- one sentence that answers the question. Say it first. If you
  remember nothing else, say this.
\item \textsf{\color{qcol}Explain} -- two or three short sentences that you can say out loud.
\item \textsf{\color{qcol}Numbers} -- only the numbers you must know by heart.
\item \textsf{\color{qcol}Derivation} -- the formula, built line by line. On the right of each line
  is the reason for the step. Write the lines on the board in this order.
\item \textsf{\color{remcol}Summary} -- a short line of key words to memorise.
\end{itemize}

\begin{tipbox}
\textsf{Five rules for the day}
\begin{enumerate}
\item Say the core sentence first, then stop for a moment. Let the examiner ask for more.
\item Say the idea in words before you write the formula.
\item Write the simplest formula first. The examiner will move you to the one he wants.
\item End an answer with the next topic: ``\ldots\ and this is why we need regularization.''
\item If you do not know, say what you know that is close: ``I do not remember the constant, but
  the idea is \ldots''
\end{enumerate}
\end{tipbox}

\textsf{Sources.} Part A: the thesis of 10 October 2026 (\texttt{main.pdf}). Part B: the JKU script
\emph{Deep Learning and Neural Networks} (DLNN, WS24, and the joint script SS25), the LSTM lecture,
and the list of past questions. Part C: the lecture \emph{Numerical Methods in Fluid Mechanics} and
its rCFD part (Prof.\ Pirker), and the figures of the thesis.
